\documentclass[hyperref=unicode,graphics=pdflatex,13pt,xcolor={usenames,dvipsnames}]{beamer}


\usepackage[usenames,dvipsnames]{xcolor}
\usepackage[noend]{algorithmic}
\usepackage{algorithm}
\usepackage{xspace}
\usepackage[algo2e,ruled,vlined,linesnumbered]{algorithm2e}
\let\oldnl\nl% Store \nl in \oldnl
\newcommand{\nonl}{\renewcommand{\nl}{\let\nl\oldnl}}% Remove line number for one line

\usepackage{amsmath}
\usepackage{mathtools}
\usepackage{cite}
\usepackage{array}
\usepackage{multirow}
\usepackage{tabularx}
\usepackage{wasysym} %smiles

\usepackage{pgfplots}
\usepgfplotslibrary{statistics}
\usetikzlibrary{calc}




\mode<presentation>
{
  % \usetheme{Goettingen}
%  \usecolortheme{crane}
  \setbeamercovered{invisible}
%  \useoutertheme{infolines}
%  \setbeamertemplate{headline}{} % removes the headline that infolines inserts
}

\graphicspath{{../pic/},{pic/}}

\setbeamertemplate{frametitle continuation}{(\insertcontinuationcount)}
\setbeamertemplate{footline}[text line]
{
  % \leavevmode%
  \hbox{%
    % \hspace{3mm}
    \begin{beamercolorbox}[ht=2.25ex,dp=1ex,left,wd=0.18\textwidth]{author in head/foot}%
        \usebeamerfont{author in head/foot}
        \insertshortauthor
    \end{beamercolorbox}%    
    \begin{beamercolorbox}[ht=2.25ex,dp=1ex,left,wd=0.63\textwidth]{title in head/foot}%
        \usebeamerfont{title in head/foot}
        \insertshorttitle
    \end{beamercolorbox}%
    \begin{beamercolorbox}[ht=2.25ex,dp=1ex,right,wd=0.2\textwidth]{framenumber}%
        \insertframenumber{}
        \hspace*{2ex} 
    \end{beamercolorbox}
  }%
}
\setbeamertemplate{navigation symbols}{}
\setbeamertemplate{itemize/enumerate body begin}{\normalsize}
\setbeamertemplate{itemize/enumerate subbody begin}{\normalsize}
\setbeamersize{text margin left=5mm,text margin right=5mm} 
% \setbeamerfont{page number in head/foot}{size=\large}
% \beamertemplatenavigationsymbolsempty

\usepackage[utf8]{inputenc}
\usepackage[T2A]{fontenc}
\usepackage[russian]{babel}
\usepackage{colortbl}
\usepackage{multirow}

\usepackage{algorithm}
\usepackage{algorithmic}

\newcommand\highlight[1]{{\color{blue}{#1}}}
\newcommand\red[1]{{\color{red}{#1}}}

\deftranslation[to=russian]{Theorem}{Теорема}

\title[Вероятностные модели, условные вероятности и независимость]{Вероятностные модели, условные вероятности и независимость}

\author[Антипов Д. С.]
{Антипов Денис Сергеевич}
\institute{
Университет ИТМО, Санкт-Петербург, Россия\\
~\\
}
\date[]{\includegraphics[height=2cm]{itmo-white-rus.png} \\
11 февраля 2021 г.}


\begin{document}

\begin{frame}[noframenumbering,plain]
%   \advance\textwidth2cm
% \hsize\textwidth
% \columnwidth\textwidth
  \titlepage
\end{frame}


\begin{frame}
  \frametitle{О лекторе}
  Обращаться --- по имени, можно на ``ты''
  \begin{itemize}
    \item PhD in computer science (Университет ИТМО и {\'E}cole Polytechnique)
    \item Занимаюсь теорией \highlight{эволюционных вычислений}
    \item Раньше вел матан, курс по теорверу --- впервые
    \item Контакты:
    \begin{itemize}
      \item telegram: \href{https://t.me/antipovden}{\highlight{\texttt{\underline{@antipovden}}}}
      \item email: \href{mailto:antipovden@yandex.ru}{\highlight{\texttt{\underline{antipovden@yandex.ru}}}}
    \end{itemize}
  \end{itemize}
\end{frame}

\begin{frame}
  \frametitle{О курсе}
  \begin{itemize}
    \item За основу взят \href{https://ocw.mit.edu/resources/res-6-012-introduction-to-probability-spring-2018/index.htm}{\highlight{\underline{курс MIT}}}
    \item Во многом схож с курсом, который преподавала раньше Ирина Александровна Суслина
    \item Немного повторяет то, что вы проходили на дискретке
    \begin{itemize}
      \item \red{Осторожно!} Можно случайно не заметить, когда началось что-то совсем для вас новое
    \end{itemize} 
  \end{itemize}
\end{frame}

\begin{frame}
  \frametitle{План лекции}
  \begin{itemize}
    \item Определение вероятностного пространства (с примерами)
    \item Условная вероятность
    \item Формула полной вероятности и формула Байеса
    \item 
  \end{itemize}

  

\end{frame}



\end{document}
